% Agent 07: Book pages 369–372 (PDF pp.13–16)
% Section 2.5 (end): Euler equation for perfect fluid, Bernoulli equation, Newtonian limit
% Section 2.6 (beginning): Radiative Transfer — notation, specific intensity,
%   mean intensity, energy density, flux, invariance of I_nu/nu^3
% Equations (2.5.30) through (2.5.42), then (2.6.1) through (2.6.11)

\documentclass[aps,prd,twocolumn,superscriptaddress,nofootinbib]{revtex4-2}
\usepackage{amsmath,amssymb,amsfonts}
\usepackage{graphicx}
\usepackage{bm}
\usepackage{mathrsfs}
\usepackage{hyperref}
\usepackage{physics}

% Custom commands
\newcommand{\fvec}[1]{\mathbf{#1}}

\begin{document}

%%% ============================================================
%%% Book page 369 — Euler equation, Bernoulli equation, Newtonian limit
%%% (PDF p.14, right column)
%%% ============================================================

\textit{Euler equation for a perfect fluid}

Consider a perfect fluid, flowing adiabatically through spacetime ($\zeta = \eta = \fvec{B} = \fvec{q} = 0$).
In this case the Euler equations $\fvec{h}\cdot(\nabla\cdot\fvec{T}) = 0$ reduce to

\begin{equation}
(\rho + p)\,\fvec{a} = -\fvec{h}\cdot\nabla p\,.
\tag{2.5.30}
\end{equation}

In words:
\begin{equation*}
\left(\!\!
\begin{array}{c}
\text{inertial mass per}\\
\text{unit volume}
\end{array}
\!\!\right)
\times\text{(4-acceleration)}
= -\left(\!\!
\begin{array}{c}
\text{pressure gradient}\\
\text{projected orthogonal to }\fvec{u}
\end{array}
\!\!\right).
\end{equation*}

\bigskip
\textit{Bernoulli equation}

Consider a perfect fluid undergoing stationary, adiabatic flow in a stationary
spacetime.  More particularly, set $\zeta = \eta = \fvec{B} = \fvec{q} = 0$; assume that spacetime is
endowed with a Killing vector field $\boldsymbol{\xi}$,

\begin{equation}
\xi_{\alpha;\beta} + \xi_{\beta;\alpha} = 0\,,
\tag{2.5.31}
\end{equation}

which need not be timelike; and assume that the flow is adiabatic, $ds/d\tau = 0$,
and stationary in the sense that

\begin{equation}
\mathscr{L}_{\boldsymbol{\xi}}\,p = 0, \quad
\nabla_{\boldsymbol{\xi}}\,\rho = \nabla_{\boldsymbol{\xi}}\,p = \cdots = 0\,.
\tag{2.5.32}
\end{equation}

Here $\mathscr{L}_{\boldsymbol{\xi}}$ is the Lie derivative along $\boldsymbol{\xi}$.  In this case the Euler equations~(2.5.30),
together with the first law of thermodynamics~(2.5.12), imply the relativistic
Bernoulli equation

\begin{equation}
d(\mu\,\fvec{u}\cdot\boldsymbol{\xi})/d\tau = 0\,;
\tag{2.5.33}
\end{equation}

i.e., $\mu\,\fvec{u}\cdot\boldsymbol{\xi}$ is constant along flow lines.

\bigskip
\textit{Newtonian limit}

The Newtonian limit of relativistic hydrodynamics is obtained when, in a nearly
global Lorentz frame, the following approximations hold:

\begin{gather}
g_{00} = -(1 + 2\Phi), \quad |\Phi| \ll 1, \quad
B^2/\rho_0 \ll 1,
\notag\\
p/\rho_0 \ll 1, \quad \Pi \ll 1, \quad v^2 \ll 1\,.
\tag{2.5.34}
\end{gather}

Here $\Phi$ is the Newtonian gravitational potential with sign $\Phi < 0$, and $v$ is the
ordinary velocity of the fluid

\begin{equation}
\fvec{v} \equiv v_j = (u^j/u^0)\,\fvec{e}_j\,.
\tag{2.5.35}
\end{equation}

In the Newtonian limit the chemical potential $\mu$ reduces to

\begin{equation}
\mu = m_B(1 + w)\,,
\tag{2.5.36}
\end{equation}

%%% ============================================================
%%% Book page 370
%%% (PDF p.15, left column)
%%% ============================================================

where $w$ is the \textit{enthalpy}

\begin{equation}
w = \Pi + p/\rho_0\,.
\tag{2.5.37}
\end{equation}

The Bernoulli equation~(2.5.33) reduces to the familiar form

\begin{equation}
\Phi + \tfrac{1}{2}v^2 + w = \text{constant along flow lines}\,.
\tag{2.5.38}
\end{equation}

The Euler equation for a perfect fluid in adiabatic flow, (2.5.30), reduces to

\begin{equation}
\frac{d\fvec{v}}{dt} = -\nabla\Phi - \frac{1}{\rho_0}\nabla p\,,
\tag{2.5.39}
\end{equation}

where $d/dt$, the derivative with respect to proper time along the flow lines, has
the Newtonian form

\begin{equation}
d/dt = \partial/\partial t + \fvec{v}\cdot\boldsymbol{\nabla}\,.
\tag{2.5.40}
\end{equation}

The first law of thermodynamics, (2.5.12) and (2.5.12$'$), reduces to

\begin{equation}
d\Pi = -p\,dV_0 + T\,ds_0\,,
\tag{2.5.41}
\end{equation}

\begin{equation}
dw = V_0\,dp + T\,ds_0\,.
\tag{2.5.42}
\end{equation}

(In Newtonian theory one usually adopts a per-unit-mass viewpoint rather than a
per-baryon viewpoint; and thus one uses $\rho_0$, $V_0$ rather than $n$, $\mathcal{V}$, $s_0$.)

%%% ============================================================
%%% Section 2.6: Radiative Transfer
%%% Book page 370–371 (PDF p.15, left–right columns)
%%% ============================================================

\subsection{Radiative Transfer}
\label{sec:2.6}

\textit{Basic references}\quad
Chandrasekhar (1960); Appendix~1 of Pacholczyk (1970); Mihalas (1970).

\bigskip
\textit{Notation and terminology}

At an arbitrary event in spacetime pick an arbitrary
local Lorentz frame.  In that frame pick an arbitrary spatial direction $\fvec{n}$ (a
unit vector; see Figure~2.6.1).  Examine the amount of energy $dE$ that is (i) carried
by photons across a unit surface area $dA$ orthogonal to $\fvec{n}$ (the surface $\mathscr{S}_m$ of
Figure~2.6.1) during unit time $dt$, with (ii) the photons having frequencies $\nu$ in
the range $d\nu$, and (iii) the photons being directed into a solid angle $d\Omega$ about
$\fvec{n}$.  The ratio

\begin{equation}
I_\nu = \frac{dE}{dt\;dA\;d\nu\;d\Omega}
\tag{2.6.1}
\end{equation}

is called the \textit{specific intensity} or simply the \textit{intensity}.  It depends on (a) the
event in spacetime; (b) the choice of Lorentz frame; (c) the direction $\fvec{n}$; (d)
the frequency $\nu$.  One also defines the \textit{total intensity} $I$ by

\begin{equation}
I = \int I_\nu\;d\nu = \frac{dE}{dt\;dA\;d\Omega}\,;
\tag{2.6.2}
\end{equation}

the \textit{average specific intensity} or \textit{average intensity} $J_\nu$ (averaged over all directions)
by

\begin{equation}
J_\nu = \frac{1}{4\pi}\int I_\nu\;d\Omega\,;
\tag{2.6.3}
\end{equation}

and the \textit{average total intensity} $J$ by

\begin{equation}
J = \frac{1}{4\pi}\int I\;d\Omega = \frac{1}{4\pi}\int I_\nu\;d\nu\;d\Omega\,.
\tag{2.6.4}
\end{equation}

It is easy to see that the \textit{energy density per unit frequency} in the radiation is

\begin{equation}
\rho^{(\text{rad})} = \frac{dE}{d^3x\;d\nu} = \frac{4\pi J_\nu}{c}\,,
\tag{2.6.5}
\end{equation}

and that the \textit{total energy density} in the radiation is

\begin{equation}
\rho^{(\text{rad})} = \frac{dE}{d^3x} = \frac{4\pi J}{c}\,.
\tag{2.6.6}
\end{equation}

%%% ============================================================
%%% Book page 371 (right column) — Figure 2.6.1, flux, invariance
%%% ============================================================

Pick a 2-surface $\mathscr{S}_m$ in the chosen Lorentz frame, with unit normal $\fvec{m}$ (Fig.\
2.6.1).  Examine the total energy $dE$ per unit area $dA$ that is carried across the
chosen surface during unit time $dt$, by photons having frequency $\nu$ in the range
$d\nu$.  Make no restrictions on the photon direction or solid angle; but count
negatively those photons that cross $dA$ from the front side toward the back.
The ratio

\begin{equation}
F_\nu = \frac{dE}{dt\;dA\;d\nu}
\tag{2.6.7}
\end{equation}

is called the \textit{specific flux}, or simply the \textit{flux}.  It is easy to see from Figure~2.6.1
that

\begin{equation}
F_\nu = \int I_\nu\,\cos\theta\;d\Omega\,.
\tag{2.6.8}
\end{equation}

\begin{figure}[htbp]
\centering
\fbox{\parbox{0.9\columnwidth}{\centering [Figure 2.6.1 placeholder]\\
The surfaces, directions, and solid angle used in the definition of intensity
and of flux.  $\mathscr{S}_m$ is a surface with unit normal $\fvec{m}$;
$\fvec{n}$ is a unit vector in the photon propagation direction;
$\theta$ is the angle between $\fvec{m}$ and $\fvec{n}$; $d\Omega$ is the solid angle element.}}
\caption{The surfaces, directions, and solid angle used in the definition of intensity
and of flux.}
\label{fig:2.6.1}
\end{figure}

The specific flux depends on (a) the chosen event in spacetime; (b) the chosen
Lorentz frame at that event; (c) the chosen 2-surface $\mathscr{S}_m$ in that Lorentz frame.

%%% ============================================================
%%% Book page 372 — Invariance of I_nu/nu^3, equation of radiative transfer
%%% (PDF p.16, left column)
%%% ============================================================

The integral of the flux over all frequencies is called the \textit{total flux}

\begin{equation}
F = \int F_\nu\;d\nu = \int I\,\cos\theta\;d\Omega\,.
\tag{2.6.9}
\end{equation}

Notice that the total flux is the ``first moment'' of the total intensity (the
``moment'' being taken with respect to the unit normal $\fvec{m}$).  One can readily
verify that the ``second moment'' is equal to the radiation pressure that acts
across the surface $\mathscr{S}_m$:

\begin{equation}
T^{(\text{rad})} \equiv \fvec{m}\cdot\fvec{T}^{(\text{rad})}\cdot\fvec{m}
= \int I\,\cos^2\theta\;d\Omega\,.
\tag{2.6.10}
\end{equation}

\bigskip
\textit{Invariance of $I_\nu/\nu^3$}

Choose a particular photon that passes through a given event in spacetime.
Observe that particular photon, and all others in the vicinity, from several different
Lorentz frames at the given event.  The frequency $\nu$ of the photon will depend
on the choice of Lorentz frame (Doppler shift from one frame to another).
The specific intensity $I_\nu$ in the frame neighborhood of the photon will also depend
on the choice of Lorentz frame.  However, the ratio $I_\nu/\nu^3$ will be Lorentz-
invariant. (Aside from a factor $h^{-4}$, $I_\nu/\nu^3$ is the invariant number density in
phase space

\begin{equation}
\mathscr{N} = \frac{dN}{d^3x\,d^3p} = \frac{1}{h^3}\,\frac{I_\nu}{\nu^3}\,;
\tag{2.6.11}
\end{equation}

see, e.g., \S22.6 of MTW.)

When photons propagate freely through curved spacetime (no interaction with
matter), the ratio $I_\nu/\nu^3$ is conserved along the world line of each photon
(Liouville's theorem).

\bibliography{agent_07}

\end{document}
